\section{Approach}
\label{sec:approach}

Figure~\ref{fig:pipeline} gives an overview. For every query, 12 signals relate the seed to a
candidate; a linear score combines them, and a \textbf{Genetic Algorithm} (GA) learns the weights of that score
from the queries of the training repositories. The same pipeline, applied to a developer's local
repository, powers the \tool{} command-line tool.

\begin{figure}[ht]
\centering
\begin{tikzpicture}[
  box/.style={draw, rounded corners, align=center, font=\footnotesize, minimum height=9mm,
              inner sep=3pt},
  arr/.style={-{Stealth[length=2mm]}, thick},
  node distance=4mm]
\node[box] (sel) {Repository\\selection\\(TerraDS)};
\node[box, right=of sel] (mine) {History mining\\(chronological,\\no future data)};
\node[box, right=of mine] (feat) {Queries +\\12 signals\\per candidate};
\node[box, right=of feat] (ga) {Genetic\\algorithm\\(learns weights)};
\node[box, right=of ga] (rec) {\tool{} \texttt{recommend}\\on a local\\repository};
\draw[arr] (sel) -- (mine);
\draw[arr] (mine) -- (feat);
\draw[arr] (feat) -- (ga);
\draw[arr] (ga) -- (rec);
\end{tikzpicture}
\caption{\small Overview: from repository selection to the recommendation tool.}
\label{fig:pipeline}
\end{figure}

\subsection{Signals}
Table~\ref{tab:features} lists the signals, all scaled to $[0, 1]$, in 3 families. The
\textbf{History} family mines \textbf{association rules} from the repository's own earlier commits, as
ROSE~\cite{zimmermann2004rose} does, at 3 granularities: between blocks, between files, and
between block types; the type-level rule (``adding a bucket goes with adding a bucket policy'') is
the only history signal available for new blocks.

The \textbf{Graph} family uses the reference graph of the repository before the commit, including the links a module call creates to another directory; shared neighbors are weighted with the Adamic--Adar index~\cite{adamic2003friends}, so a shared rare neighbor counts more than a shared hub such as a common tags variable. The propagation rule is mined: it learns from earlier commits how often a direct neighbor changes, separately for added, modified and removed seeds and for each edge direction, because a removed block breaks its referencers far more often than an edited one.

The \textbf{Text} family compares names and the contents of the blocks (TF-IDF over the tokens of their bodies), which links copies of a configuration that no reference connects. Every signal is computed per repository and only from data available before the commit (Section~\ref{sec:problem}).

\begin{table}[ht]
\begin{tabular}{@{}l>{\raggedright\arraybackslash}p{10.2cm}@{}}
\toprule
Signal & Definition \\
\midrule
\multicolumn{2}{@{}l}{\emph{History} (earlier commits of the repository)} \\
\textbf{co-change confidence} & share of past commits changing $s$ that also changed $c$ \\
\textbf{co-change count}      & number of past commits changing both (log-scaled) \\
\textbf{co-change recency}    & co-changes weighted by how recent they are \\
\textbf{recent changes}      & how often $c$ changed in the last 200 commits \\
\textbf{file confidence}      & the same rule between the files of $s$ and $c$ \\
\textbf{type rule}            & rule between block types, e.g.\ modified \texttt{variable} $\Rightarrow$ changed \texttt{module}; for types to add, $\Rightarrow$ added type \\
\midrule
\multicolumn{2}{@{}l}{\emph{Graph} (references before the commit)} \\
proximity            & $1/(1+d)$ for reference distance $d \le 4$, else 0 \\
\textbf{shared neighbors}    & Adamic--Adar score of common neighbors \\
\textbf{degree}              & number of references of $c$ \\
\textbf{propagation rule}     & learned rate at which direct neighbors change, per seed action and edge direction \\
\midrule
\multicolumn{2}{@{}l}{\emph{Text}} \\
\textbf{name similarity}      & overlap of name tokens; for types to add, seed name vs.\ type name \\
\textbf{content similarity}   & TF-IDF cosine of the seed's new body and the candidate's body \\
\bottomrule
\end{tabular}
\centering
\footnotesize
\caption{The 12 engineered features relating a seed $s$ to a candidate $c$.}
\label{tab:features}
\end{table}

\subsection{Scoring Model}
A candidate with signal vector $\mathbf{x}$ receives the score
\[
  \mathit{score}(\mathbf{x}) = \mathbf{w}^\top \mathbf{x} + b \cdot [\text{$c$ is a type to add}],
  \qquad w_i \ge 0,\ \textstyle\sum_i w_i = 1,\ b \in [-1, 1].
\]
Each weight reads as the share of trust in one signal. The offset $b$ is needed because a type to
add has only two non-zero signals (type rule and name similarity), so without it such candidates
would always rank below existing blocks, or above them if those two weights grew. Weights are
non-negative because every signal was designed so that a higher value means a stronger relation;
a check that allows negative weights confirms this (Section~\ref{sec:results}).

\subsection{Genetic Algorithm}
The 13 parameters ($\mathbf{w}$ and $b$) are learned by a Genetic Algorithm implemented from
scratch. A chromosome holds the 13 real-valued genes; after every variation step the
weights are clipped at zero and normalized to sum 1, and the offset is clipped to $[-1, 1]$. The
fitness is the \emph{MRR} obtained with the chromosome's weights on a fixed sample of 5,000 training
queries, with the same number of queries from every training repository so that large
repositories do not dominate.

Evaluating on a sample instead of all ${\sim}44$k training queries is
a form of fitness approximation~\cite{jin2005fitness}: a check on one fold gave the same test MRR
at a ninth of the cost. The fitness of the whole population is computed with one matrix product.
Parents are chosen by tournament selection (size 3) and recombined with BLX-$\alpha$ crossover
($\alpha = 0.5$)~\cite{eshelman1993blx}; Gaussian mutation uses a step size that shrinks from 0.2
to 0.02 over the run, and the two best individuals survive unchanged (elitism).

The remaining settings were not assumed but tuned on training data only. A convergence check with
a generous budget (200 generations) showed that validation MRR barely improves after 60 generations
of 60 individuals (0.735 vs.\ 0.738), which fixed the budget at about 3,600 fitness evaluations. Within each \textbf{cross-validation fold}, a \textbf{grid search} over population/generations at that budget, mutation rate and crossover rate (27 configurations, 3 runs each) trained on 75\% of the training repositories and selected the
configuration with the best MRR on the other 25\%. The final model, trained on all 53 repositories,
uses a population of 30, 120 generations, a mutation rate of 0.1 and a crossover rate of 0.8.

\subsection{The \tool{} Tool}
The command \texttt{tfimpact recommend --repo PATH --diff} applies the final model to a local
Terraform repository. It mines the repository's history with the same pipeline, takes the blocks
changed in the working tree as seeds (or a single block given with \texttt{--block}), and prints
the top blocks to modify and block types to add; \texttt{--explain} shows the signals behind each
suggestion. It needs only Git and the 13 learned numbers, and runs in a few seconds on a
repository with a few hundred commits.
